\chapter{Discussion and Analysis}
\label{ch:discussion}

% ============================================================================
\section{Why Gating Always Hurts: A Structural Argument}
\label{sec:gating_argument}
% ============================================================================

The AND-fusion penalty (Eq.~\ref{eq:and_penalty}) is not an artifact of specific
sensors or databases. It is a structural consequence of treating a fallible
biometric as a mandatory gate. Let $G_\mathrm{FP}$ and $G_\mathrm{Iris}$
be independent Bernoulli random variables indicating genuine acceptance. Under
AND-fusion:
\begin{equation}
  \Pr[\mathrm{Accept}] = G_\mathrm{FP} \wedge G_\mathrm{Iris}
  \leq G_\mathrm{FP}
\end{equation}
with equality only if $G_\mathrm{Iris} = 1$ always, which is impossible for
any realistic biometric. The tighter the coupling (Architectures B, C, D), the
more strictly this equality becomes an inequality.

Architectures B--D impose a stronger bound: their EER is floor-bounded by the
iris FRR. At $B_s = 255$, iris FRR $\approx 38\%$; therefore:
\begin{equation}
  \mathrm{EER}_\mathrm{B/C/D} \geq 0.38 \times (1 - \mathrm{FAR}_\mathrm{FP})
  \approx 38\%
\end{equation}
for any fingerprint system, regardless of sensor quality. The experimental results
(EER $\geq 19\%$ under the best fingerprint conditions) closely track this bound.

IBFV breaks the structural constraint by making the iris conditional: the iris
path is activated only when the iris key recovers successfully. This converts an
AND-gate ($\text{must pass both}$) into a conditional booster ($\text{may use
iris if available}$).

% ============================================================================
\section{Where Boosting Helps Most}
\label{sec:boost_analysis}
% ============================================================================

The 61.3\% relative EER reduction on FVC2002~DB3 (vs.\ $\sim$40--60\% on
other databases) is explained by the interaction between fingerprint difficulty
and iris bonus:

\begin{itemize}
  \item On an \textbf{easy database} (DB1, unimodal EER already 0.000\%), there
  is no room for further absolute improvement. Both unimodal and IBFV achieve
  0.000\% EER; the zero-penalty guarantee holds trivially, and bonus points do
  not elevate impostor FAR on these high-quality databases.
  \item On a \textbf{challenging database} (DB3, unimodal EER~2.762\%), many genuine
  queries fail because they recover too few matching fingerprint points to
  enumerate $k{+}1$ subsets successfully. Adding $K = 4$ bonus points reduces the
  required fingerprint match from $k{+}1 = 11$ to $k{+}1{-}K = 7$, which is a
  qualitatively easier threshold. This corresponds to a $\Delta \approx 4$-point
  shift in the distribution of matched points, moving substantially more genuines
  above the acceptance threshold.
  \item \textbf{Diminishing returns} at $K > 6$: once $K$ reaches $k{+}1$, all
  required polynomial points can be supplied by the iris alone. Further increasing
  $K$ cannot reduce the subset enumeration further; the fingerprint constraint
  is already fully satisfied.
\end{itemize}

% ============================================================================
\section{Security Analysis}
\label{sec:security_analysis}
% ============================================================================

\subsection{Irreversibility}
\label{sec:irreversibility}

The ECC vault protects minutiae via ECDLP: recovering quantized minutia $q_i$
from stored $x$-coordinate $P_i.x$ requires solving ECDLP on Brainpool~P256r1.
The best known attack requires $\approx 2^{128}$ group operations---beyond the
computational reach of any foreseeable adversary.

The iris commitment stores only helper data $\mathbf{h}$ and key hash $H(\mathbf{k})$.
The key space for $n_b = 514$ blocks is $2^{514}$; brute-forcing the iris key
is infeasible. Together, the vault and iris commitment achieve the irreversibility
requirement of ISO/IEC~24745 \cite{iso24745}.

\subsection{Unlinkability}
\label{sec:unlinkability}

Unlinkability is quantified by Meng et al.'s $D_{\leftrightarrow}$ metric and the
$D_\mathrm{sys}$ system-level measure \cite{gomezbarrero2017}. A threshold of
$D_\mathrm{sys} \leq 0.05$ is required for unlinkability.

The IBFV vault achieves $D_\mathrm{sys} \leq 0.03$ through three mechanisms:
\begin{itemize}
  \item \textbf{EC mapping}: Vault x-coordinates are elliptic curve points; no
  two vaults share the same x-coordinates (fresh polynomial each enrollment).
  \item \textbf{Virtual x-coordinate randomness}: Bonus point coordinates depend
  on $H(\mathbf{k})$; with a fresh iris commitment key per enrollment, they are
  computationally independent across templates.
  \item \textbf{Chaff randomness}: Fresh random chaff per enrollment ensures no
  statistical relationship between vaults.
\end{itemize}

\subsection{Revocability}
\label{sec:revocability}

Re-enrollment with new quantization factor $f'$, new polynomial, and new iris
commitment key produces a vault in a different domain. By the security of SHA-256
(used in key hashing) and by Theorem~2 (Shamir share confidentiality), new and
revoked shares are information-theoretically independent. Any biometric data that
an adversary extracted from the old vault provides zero advantage for attacking
the new vault.

\subsection{D-IBFV Security Summary}
\label{sec:dibfv_security_summary}

\begin{table}[h]
\centering
\caption{D-IBFV security properties per ISO/IEC~24745.}
\label{tab:security_summary}
\small
\begin{tabular}{@{}p{3.2cm}p{8.2cm}p{3.8cm}@{}}
\toprule
\textbf{Property} & \textbf{Mechanism} & \textbf{Quantification} \\
\midrule
Irreversibility     & ECDLP + iris key space & $2^{128}$ / $2^{514}$ \\
Unlinkability       & Vault randomness + IBFV PRF & $D_\mathrm{sys} \leq 0.03$ \\
Revocability        & Re-enrollment protocol & Zero information leakage (Thm.~2) \\
Brute-force defense & Expected $6.8{\times}10^{8}$ attempts at $10^6$/s $\approx 680$\,s; mitigations: HSM, $k{=}15$, Argon2 & Known limitation~\cite{tams2013attacks} \\
Server compromise   & Shamir $(t,n)$: $<t$ shares give nothing & Thm.~2: info-theoretic \\
IPFS tampering      & CID integrity + SHA-256 trailer & Detect any 1-bit change \\
Blockchain tamper   & Consensus + append-only & Raft/BFT: $< 50\%$ attack budget \\
\bottomrule
\end{tabular}
\end{table}

\noindent\textbf{Brute-force resistance.} An adversary holding a complete vault
must identify $k{+}1$ genuine points among $|\mathcal{V}| = 6g + K = 124$ total
points. The expected number of attempts is
$\binom{124}{11}/\binom{24}{11} \approx 6.8 \times 10^{8}$; at $10^6$ attempts
per second this requires $\sim\!680$~seconds---\emph{feasible} for an offline
attacker, a well-known limitation of all fuzzy vault schemes~\cite{tams2013attacks}.
Standard mitigations include: (1)~storing $S$ in a hardware security module (HSM)
so vault compromise alone is insufficient; (2)~increasing polynomial degree $k$
(at $k = 15$, cost rises to $\sim\!3 \times 10^{6}$ years); and
(3)~adding key-stretching (e.g., Argon2). The primary template protection
goal---preventing \emph{biometric recovery}---is achieved regardless, since
reconstructing $p(x)$ reveals only the secret $s$, not the original minutiae
(ECDLP on Brainpool~P256r1 requires $\sim\!2^{128}$ operations).

% ============================================================================
\section{Comparison with Prior Systems}
\label{sec:prior_comparison}
% ============================================================================

\begin{table}[h]
\centering
\caption{Comparison with key prior multimodal fuzzy vault and blockchain biometric systems.}
\label{tab:prior_comparison}
\small
\begin{tabular}{@{}lccccc@{}}
\toprule
\textbf{System} & \textbf{Zero-Penalty} & \textbf{Decentralized} & \textbf{Formal Proofs}
 & \textbf{Best EER} & \textbf{ISO 24745} \\
\midrule
Nandakumar \cite{nandakumar2008multibiometric} & No & No & No & 2.14\% & Partial \\
Moon \cite{moon2010multibiometric}             & No & No & No & 4.50\% & Partial \\
Kelkboom \cite{kelkboom2011multialgorithm}               & No & No & No & 3.80\% & Partial \\
Maurya \cite{maurya2025enhancing}              & N/A & No & No & 0.005\% & Yes \\
Delgado-Mohatar \cite{delgado2020blockchain}   & N/A & Yes & No & --- & No \\
Sharma 2024 \cite{sharma2024multimodal}        & N/A & Yes & No & --- & Partial \\
\midrule
\textbf{D-IBFV (proposed)} & \textbf{Yes} & \textbf{Yes} & \textbf{Yes} &
  \textbf{0.000\%} & \textbf{Full} \\
\bottomrule
\end{tabular}
\end{table}

D-IBFV is the only system in the comparison that simultaneously: (1) guarantees
zero-penalty multimodal fusion, (2) operates on a decentralized blockchain
infrastructure with Shamir threshold sharing, and (3) provides formal
information-theoretic proofs.

% ============================================================================
\section{Limitations and Caveats}
\label{sec:limitations}
% ============================================================================

\begin{enumerate}
  \item \textbf{Chimeric database}: Fingerprint and iris subjects are paired
  artificially rather than co-acquired from the same individuals. The independence
  assumption may not hold for all real deployments where inter-modal correlation
  could affect security estimates. The fusion is evaluated on balanced impostor
  scenarios, which may differ from real-world attack distributions.

  \item \textbf{Iris GAR ($\sim$62\%)}: The iris fuzzy commitment achieves only
  $\sim$62\% genuine acceptance at $B_s = 255$. While the zero-penalty guarantee
  eliminates any \emph{penalty} from this limitation, higher iris GAR would
  translate to a larger IBFV boost. The iris pipeline does not include
  iris-specific quality filters (e.g., specular reflection removal, pupil dilation
  normalization) that could improve GAR.

  \item \textbf{Gas cost at Ethereum L1}: At current prices ($\sim$\$18.80 per
  enrollment), Ethereum L1 deployment is cost-prohibitive for large-scale systems.
  Layer-2 solutions (Polygon PoS, Arbitrum) reduce costs by $\sim 10{,}000\times$
  to sub-dollar, but L2 security assumptions differ from L1 finality guarantees.

  \item \textbf{Ganache vs.\ mainnet}: Ganache provides a local single-node
  EVM simulation. Mainnet deployment would introduce additional latency from
  network propagation, mempool congestion, and variable gas prices.

  \item \textbf{Sample size}: The biometric evaluation uses 100 subjects per
  FVC database and 249 subjects for CASIA. Larger-scale evaluations (thousands
  of subjects) may reveal distribution tail effects not visible at these scales.
\end{enumerate}


